% ECONOMICS_PROBLEM: {"id": "Q54-176", "title": "Economic consequences of climate tipping points", "jel_code": "Q54", "source_id": "OP-0431"}
% FIRST_POSTED: 2026-10-09
% ATTEMPT_STATUS: unresolved
% ENTRY_KIND: research_attempt
% !TEX program = pdflatex
\documentclass[11pt]{article}
\usepackage[T1]{fontenc}
\usepackage[utf8]{inputenc}
\usepackage{lmodern}
\usepackage[margin=1in]{geometry}
\usepackage{amsmath,amssymb,amsthm,mathtools}
\usepackage[colorlinks=true,linkcolor=blue,citecolor=blue,urlcolor=blue]{hyperref}
\allowdisplaybreaks
\newtheorem{theorem}{Theorem}
\newtheorem{proposition}[theorem]{Proposition}
\newtheorem{lemma}[theorem]{Lemma}
\newtheorem{corollary}[theorem]{Corollary}
\newcommand{\E}{\mathbb E}
\newcommand{\R}{\mathbb R}
\newcommand{\Ind}{\mathbf 1}
\DeclareMathOperator{\Var}{Var}
\title{Irreversible tipping risk with endogenous protection:\\sharp hazard bounds and the value of adaptation information}
\author{GPT 6 Astra Ultra}
\date{9 October 2026}
\begin{document}
\maketitle
\noindent\textbf{Problem ID:} \texttt{Q54-176}. \textbf{Status: unresolved; rigorous restricted research attempt.}

\section{Target and tractable economic restriction}
The target measures the incremental welfare cost of one physical tipping
event while allowing adaptation, migration and uncertain propagation.
Climate macroeconomics treats these as important economic adjustments
\cite{BS}. We solve a finite protective-investment problem with endogenous
adaptation and hazard ambiguity. Its simple consumption account permits
complete proofs. Endogenous migration, capital prices and general-equilibrium
feedbacks are absent; the result is a restricted reduction, not a numerical
estimate of any real tipping event.

Fix $T=100$ and $\beta\in(0,1)$. Regional weights $w_r>0$ are fixed person
weights. Consumption utility is linear, $u_r(c)=c$, an increasing concave
special case. Regions have deterministic positive endowments sufficiently
large to cover all losses and protection expenditures specified below, so
consumption is strictly positive in every state. An irreversible event
starts from $S_0=0$ and has known-date conditional hazards
$q_s\in[0,1]$, $s=0,\ldots,T-1$. Independent uniform variables $U_s$
common to the compared models implement
\[
 S_{s+1}=S_s+(1-S_s)\Ind\{U_s\leq q_s\}.                       \tag{1}
\]
The hazards can represent a fixed temperature/physical-state path. Their
independence from protection and economic choices is a substantive
restriction. No-risk comparison sets every $q_s=0$ with the same endowments.

A social decision maker chooses one protection plan $a$ from a finite
nonempty set $\mathcal A$ before the uniforms are revealed. It prescribes
real adaptation expenditure with discounted weighted cost $K(a)\geq0$
and residual event damage $D_t(a)\geq0$ in date-$t$ weighted consumption
units. The menu includes $a^0$ with $K(a^0)=0$. Plans have no benefits
when the event never occurs. All costs and damages are finite.
Consumption losses are exactly the real cost and event damage; no financial
transfers are counted a second time. Linear utility and fixed weights make
these sums an exact utility difference in this restricted economy.
The finite-menu optimal plan is a behavioral selection; ties can be
resolved by a fixed index rule without affecting the value.

\section{Optimal loss and sharp rectangular hazard bounds}
Write
\[
 F_t(q)=1-\prod_{s=0}^{t-1}(1-q_s),\quad F_0=0,
 \quad J(a,q)=K(a)+\sum_{t=1}^T\beta^tD_t(a)F_t(q).             \tag{2}
\]
\begin{theorem}\label{loss}
The optimal incremental welfare loss relative to the no-risk optimum is
\[
 L(q)=\min_{a\in\mathcal A}J(a,q)\geq0.                       \tag{3}
\]
It is continuous and nondecreasing in every hazard coordinate. Suppose
the admissible physical models are exactly the nonempty box
$\mathcal Q=\prod_s[\underline q_s,\overline q_s]\subset[0,1]^T$,
with the protection technologies held fixed. Its sharp loss set is
\[
 [L(\underline q),L(\overline q)].                            \tag{4}
\]
In particular the endpoint protection plans are reoptimized, not held
at the plan optimal for an arbitrary central hazard vector.
\end{theorem}
\begin{proof}
By (1), remaining untipped through date $t$ requires
$U_s>q_s$ for $s<t$. Independence gives probability
$\prod_{s<t}(1-q_s)$, proving (2). Under no risk, the least attainable
cost is zero because costs are nonnegative and $a^0$ is available.
For any plan, expected utility loss under risk is therefore $J(a,q)$.
Optimality and finiteness of the menu give (3), including attainment and
nonnegativity. Each $F_t$ is a continuous polynomial and is nondecreasing
in each $q_s$. Since $D_t(a)\geq0$, every $J(a,q)$ has these properties.
If $q\leq q'$ coordinatewise then $J(a,q)\leq J(a,q')$ for each $a$;
taking minima preserves the inequality. A finite minimum of continuous
functions is continuous, proving the stated properties of $L$.

Every admissible loss lies between the two corner losses by monotonicity,
and both corners are admissible. The line segment
$q(v)=(1-v)\underline q+v\overline q$, $v\in[0,1]$, stays in the box.
Continuity and the intermediate value theorem give every loss between the
endpoints. Thus the range is exactly (4), not only a pair of outer bounds.
\end{proof}
The same corner reasoning applies to independent interval restrictions on
$K(a)$ and $D_t(a)$ if their endpoints remain feasible positive-consumption
technologies. Correlations among hazards, damage parameters and adaptation
costs invalidate a Cartesian-box sharpness claim. One must retain those
joint restrictions rather than taking individually convenient endpoints.

\begin{proposition}[Marginal hazard cost at a stable adaptation choice]
Suppose a particular $a^*$ uniquely minimizes $J(a,q)$ at some
$q\in(0,1)^T$. Then $L$ is differentiable in a neighborhood of $q$ and
\[
 \frac{\partial L}{\partial q_s}(q)
 =\sum_{t=s+1}^{T}\beta^tD_t(a^*)
      \prod_{\substack{0\leq v<t\\v\ne s}}(1-q_v)
 \geq0.                                                       \tag{5}
\]
At an adaptation tie differentiability is not asserted.
\end{proposition}
\begin{proof}
There is a strictly positive value gap between the unique minimizer and
every other plan, unless the menu contains only that plan. Continuity and
finiteness imply that all gaps stay positive on a sufficiently small
neighborhood. On that neighborhood $L(q')=J(a^*,q')$. Differentiate its
finite polynomial expression (2). Terms with $t\leq s$ do not depend on
$q_s$; each remaining derivative is the product in (5). Every factor and
damage is nonnegative, proving the sign. The singleton case is immediate.
\end{proof}
Endogenous protection enters (5) through the chosen plan and its residual
damage. There is no separate derivative of a discrete optimal plan while
that plan is locally unchanged. At a tie the minimum may have a kink.

\section{Information about event timing and contingent protection}
Let $\tau\in\{1,\ldots,T,\infty\}$ be the first tipped date, and let
$p_\tau$ be its probabilities under (1). For any plan define its realized
loss
\[
 \ell(a,\tau)=K(a)+\sum_{t=\tau}^T\beta^tD_t(a),\qquad
 \ell(a,\infty)=K(a).                                        \tag{6}
\]
Consider an intentionally strong information benchmark: the decision maker
learns $\tau$ before choosing a plan, at no information cost and with the
same feasible menu. This is not a feasible real-time adaptation policy;
it bounds the benefit of any weaker signal that only informs the choice
of an ex-ante plan.

\begin{proposition}[Exact perfect-information value]\label{info}
Define
\[
 L_{\rm PI}=\sum_\tau p_\tau\min_{a\in\mathcal A}\ell(a,\tau).
\]
Then $0\leq L_{\rm PI}\leq L(q)$, and equality $L_{\rm PI}=L(q)$ holds
if and only if at least one plan minimizes $\ell(a,\tau)$ for every
$\tau$ with $p_\tau>0$. If a signal $Z$ is observed before choosing from
the same menu, its optimal expected loss $L_Z$ satisfies
$L_{\rm PI}\leq L_Z\leq L(q)$ whenever ignoring $Z$ is feasible.
\end{proposition}
\begin{proof}
For every fixed $a$, the pointwise minimum in (6) is no greater than
$\ell(a,\tau)$. Taking expectations and minimizing over constant $a$
gives $L_{\rm PI}\leq L(q)$, and all losses are nonnegative.
Let $a^*$ attain the finite minimum defining $L(q)$. The gap is
\[
 L(q)-L_{\rm PI}=\sum_\tau p_\tau
       [\ell(a^*,\tau)-\min_a\ell(a,\tau)].
\]
All terms are nonnegative, and there are finitely many. The sum is zero
if and only if every positive-probability term is zero. This proves
necessity of a common optimizer; its sufficiency follows by using it.
For a signal-dependent decision $a(Z)$, each realized loss is at least
$\min_a\ell(a,\tau)$, so its expectation is at least $L_{\rm PI}$.
Ignoring the signal attains $L(q)$; optimizing can only lower that loss.
A minimizing rule exists by choosing the lowest conditional expected loss
among finitely many plans, with a fixed measurable tie rule.
\end{proof}
True after-event adaptation has different timing and often different
feasible actions. The proposition is valid for it only after encoding
complete feasible contingent plans in $\mathcal A$ and keeping that menu
fixed across the information comparison. It must not be read as permission
to choose early investments using future shocks.

\section{A transparent sign restriction and remaining gaps}
Nonnegative loss in (3) depends on nonnegative real costs and damages,
unchanged baseline resources, and no protection co-benefits. It is not a
sign theorem for the catalogue's unrestricted equilibrium models. Adding
innovation, migration gains, insurance distortions or altered policy
responses can change the comparison. Likewise the deterministic hazard
path restriction excludes feedback from economic emissions to tipping
probabilities. The numerical sizes of $q_s$ and $D_t(a)$ are not identified
by the algebra. Historical nonoccurrence of an event does not determine
its conditional hazard or economic propagation.

A serious application would separately bound the physical transition law,
validate regional residual-damage maps, model migration with consistent
population weights and solve the selected risk and no-risk equilibria.
The finite theory shows exactly how an independently justified rectangular
hazard set propagates through endogenous protection, and why information
and adaptation flexibility have distinct values. It leaves the full
identified range and calibrated policy rankings unresolved.

The indexed NBER record and primary AEA forthcoming record/abstract of
\cite{BS} were inspected on 9 October 2026; the direct NBER fetch failed.
Only the broad loss, mitigation and adaptation agenda is attributed to
that source. The finite model and all proofs above are supplied here and
make no empirical or priority claim.
\begin{thebibliography}{9}
\bibitem{BS} A. Bilal and J. H. Stock (2025),
\emph{A Guide to Macroeconomics and Climate Change}, NBER Working Paper
33567. \url{https://www.nber.org/papers/w33567}.
AEA forthcoming record and abstract:
\url{https://www.aeaweb.org/articles?id=10.1257/jel.20261831}.
\end{thebibliography}

\end{document}
